\documentclass[12pt]{article}
\usepackage[utf8]{inputenc}
\usepackage[T1]{fontenc}
\usepackage{hyperref}
\title{Ngumbarl Language Resources for K-12 Curriculum Development}
\author{Compiled by Claude Sonnet 5.5 (claude-sonnet-5-5) for Andr\'as Kornai}
\date{}
\begin{document}
\maketitle

\section{Introduction}
Ngumbarl (ISO 639-3 xnm; Glottocode ngum1253) is an extinct, poorly documented Nyulnyulan language of the coast north of Broome, Western Australia. It was documented in the late 1960s and had few or no speakers by 1984. Sources disagree on its branch: McGregor is cited as placing it in Western Nyulnyulan, while Bowern argues from the Daisy Bates and Billingee records that it is Eastern Nyulnyulan.

\section{Primary Educational Resources}
No primers, textbooks or school programmes were found in search results.

\section{Linguistic and Reference Materials}
\begin{itemize}
\item Ngumbarl language (Wikipedia): \url{https://en.wikipedia.org/wiki/Ngumbarl_language}
\item AIATSIS Austlang entry K4, Ngumbarl: \url{https://collection.aiatsis.gov.au/austlang/language/K4}
\item Nyulnyulan languages (family overview, Wikipedia mirror): \url{https://webcf.waybackmachine.org/web/20220216151533/https://en.wikipedia.org/wiki/Nyulnyulan_languages}
\end{itemize}
The main surviving material appears to be a word list recorded by Daisy Bates and Billingee in the early twentieth century. I did not see a direct link to it.

\section{Religious and Community Resources}
\begin{itemize}
\item Ngumbarl people (Wikipedia): \url{https://en.wikipedia.org/wiki/Ngumbarl_people}
\end{itemize}

\section{Digital and Audio Resources}
None found.

\section{Curriculum Development Context}
There are no known speakers and documentation is thin. A school programme could cover the people, their country and the language's history, and could draw on the Bates word list through archives. Any language content should be developed with Aboriginal communities of the Broome region. Nothing found supports a full language curriculum.
\end{document}